% Migrated from the Word baseline; do not duplicate content in other section files.
\section{结论}\label{sec:conclusion}

本文围绕磁驱柔性导管的多物理大变形动力学，构建了从几何精确连续模型、联合 Galerkin 约化、正则 Hamiltonian 表述到结构化 CHNN 和 KAM 稳定性分析的统一框架。该框架保留截面姿态、独立扭转和转动惯量，并通过排他能量分账区分保守作用、被动耗散、背景流与控制输入，从而为连续模型、COMSOL 数据和有限维学习系统提供一致的状态与功率口径。

复值表示增强的 CHNN 采用实值物理 Hamiltonian 读出，并以独立分支表示保守修正、耗散和外部端口。统计学习误差只是 KAM 接口的一部分；完整认证还依赖高阶或解析误差证书、既有物理摄动、坐标与截断误差、辛离散误差、非退化条件和 Diophantine 条件。由此，本文给出了连续 Hamiltonian KAM、辛映射 KAM 与后验不变环面核验之间职责明确的分析路径。

P0、P1、P1.5 和 P1.6 给出了低维或三轨迹小样本结果。P1 在 derivative-only 比较中显示条件性的 CHNN 优势；P1.5 显示 residual 相对 physics-only 的显著收益，但没有稳定的一步场优势；P1.6 则显示 \(30/30\) 个约束模型精确满足平衡条件而预测误差明显恶化。这些结果不能外推为完整 P2 任务上的普遍 CHNN 优势，且不同随机种子不等于独立物理样本。

COMSOL--Galerkin bridge 的 \(0.5\%\) 扰动能量判据并非三条轨迹均满足，而独立的 \(2\%\) force-v2 联合 L2 判据在三条轨迹上均满足，两者应保持不同解释。经过同轨迹比较、误差来源隔离、余切投影识别、求解器输出语义比较和固定基初值修正后，三层时间步计算得到 \(1.64186\)--\(2.00088\) 的观察收敛阶；在 \(1\,\mathrm{s}\) 工况上，G2、G3 点态、G3 impulse 和扰动能量峰峰漂移分别为 \(0.00146226\)、\(0.000221162\)、\(0.0000563606\) 和 \(0.00307057<0.005\)。上述验证表明，当前 canonical 数据合同满足后续多幅值实验的数据生成要求。

多幅值统计、完整 96 轨迹、P2 模型比较、开放端口、高保真偏差和数值 KAM 仍需后续数值工作。本文尚未证明 BDF 为辛积分器，尚未取得数值 KAM 认证、不变环面或工程安全域，也未证明 CHNN 在完整 P2 任务上普遍优于 Real HNN。本文的理论保证限于满足相应结构、正则性、非共振和小扰动条件的保守骨架；开放端口、接触切换、模型偏差和工程安全性需结合能量收支、端口指标、高保真验证及实验数据综合评价。
